\documentclass[12pt,a4paper]{article}
\usepackage[utf8]{inputenc}
\usepackage[T1]{fontenc}
\usepackage[brazil]{babel}
\usepackage{lmodern}
\usepackage[margin=2.5cm]{geometry}
\usepackage{microtype}
\usepackage{amsmath,amssymb}
\usepackage{booktabs,tabularx,array}
\usepackage{graphicx,xcolor}
\usepackage{listings}
\usepackage[hidelinks]{hyperref}
\setlength{\parskip}{0.35em}
\setlength{\parindent}{1.25cm}
\newcommand{\pendente}{\textemdash}
\lstset{basicstyle=\ttfamily\small,breaklines=true,frame=single,
  columns=fullflexible,keepspaces=true,showstringspaces=false}
\newcommand{\grafico}[3]{
  \begin{figure}[htbp]
    \centering
    \IfFileExists{#1}{\includegraphics[width=0.9\linewidth]{#1}}{
      \fbox{\parbox[c][3cm][c]{0.85\linewidth}{\mbox{}}}}
    \caption{#2}\label{#3}
  \end{figure}}

\begin{document}
\begin{titlepage}
\centering
{\large Universidade de São Paulo\par}
{\large Instituto de Ciências Matemáticas e de Computação\par}
\vspace{1cm}
{\large SSC0903 --- Computação de Alto Desempenho\par}
{\large Primeiro Trabalho Prático\par}
\vfill
{\LARGE\bfseries Incêndio Florestal\par}
\vspace{0.5cm}
{\Large Implementações sequencial e paralela com OpenMP\par}
\vfill
\begin{tabular}{ll}
\textbf{Integrante} & \textbf{Número USP} \\
João Ricardo de Almeida Lustosa & 15453697 \\
Catarina Moreira Lima & 8957221 \\
{Nome completo} & {Número USP} \\
\end{tabular}
\vfill
São Carlos\par
Setembro de 2026
\end{titlepage}
\tableofcontents
\newpage

\section{Introdução}

Simulações em grades discretas permitem representar a evolução de um sistema
por meio de regras locais aplicadas repetidamente a suas células. Neste
trabalho, uma matriz retangular representa uma área florestal na qual o fogo
se propaga a partir de focos iniciais. A evolução depende da cobertura e da
umidade de cada célula, da presença de vizinhos em chamas, da direção e da
intensidade do vento e da ativação de zonas de contenção em instantes
predeterminados~\cite{enunciado}.

O modelo utiliza passos discretos e atualização síncrona: o resultado de uma
célula em um passo é calculado a partir da configuração atual, sem utilizar
resultados parcialmente produzidos para outras células. Essa propriedade
permite distribuir as atualizações entre threads, desde que sejam respeitadas
as dependências entre as etapas. A ativação das contenções deve terminar
antes da propagação, e a atualização completa deve preceder a troca das
estruturas que representam as configurações atual e seguinte.

A atividade relaciona, assim, correção e desempenho em programação paralela
com memória compartilhada. A distribuição das células pode reduzir o tempo
de processamento, mas introduz custos de sincronização e escalonamento.
Além disso, o trabalho realizado por uma célula depende de seu estado:
células intactas avaliam a vizinhança, enquanto outros estados exigem
transições mais simples. A quantidade de threads, a política de distribuição
das iterações e o tamanho da grade são, portanto, aspectos relevantes para
a avaliação experimental.

O projeto contempla duas implementações em C: uma sequencial, utilizada
como referência de comparação, e outra paralela com OpenMP. Ambas devem
produzir os mesmos resultados para uma mesma entrada, exceto pelo tempo de
execução. Embora a inicialização empregue geração pseudoaleatória, a
semente e a ordem das chamadas são fixadas, e não há decisões aleatórias
durante a propagação. O modelo é o definido para a disciplina e não pretende
representar um modelo físico validado de incêndios florestais.

\section{Objetivos}

\subsection{Objetivo geral}

Implementar e avaliar uma simulação determinística da propagação de incêndio,
comparando uma solução sequencial em C com uma solução paralela em C/OpenMP,
de modo a preservar as regras do modelo e investigar o desempenho em
diferentes configurações de execução.

\subsection{Objetivos específicos}

\begin{itemize}
\item Representar a floresta em armazenamento linear e implementar a leitura,
validação e inicialização dos dados conforme o enunciado.
\item Implementar as transições de estado, a influência direcional do vento,
os tempos de queima e a ativação de zonas de contenção.
\item Desenvolver uma versão sequencial que sirva de referência para a
validação da versão paralela, complementada por testes das regras do modelo.
\item Distribuir a ativação e a atualização das células em uma região
paralela persistente, utilizando \texttt{omp for}, SIMD, reduções e
sincronização adequada.
\item Verificar a equivalência das saídas entre versões e a independência
dos resultados em relação à quantidade de threads e ao schedule.
\item Comparar pelo menos duas políticas de escalonamento e analisar tempos,
speedup e eficiência em diferentes cargas e quantidades de threads.
\end{itemize}

\section{Modelagem do problema}\label{sec:modelagem}

\subsection{Domínio espacial e temporal}

A área é representada pelo conjunto de células
\begin{equation}
\Omega=\{(l,c)\mid 0\leq l<L,\;0\leq c<C\},
\end{equation}
com $N=LC$ posições. O armazenamento é linear, com índice $i=lC+c$.
As coordenadas podem ser recuperadas pela divisão inteira $l=i/C$ e pelo
resto $c=i\bmod C$. Essa representação permite percorrer a grade por um
único laço e distribuir suas iterações entre threads.

A evolução ocorre em, no máximo, $P$ passos, numerados de $0$ a $P-1$.
Uma configuração inicial é construída antes do primeiro passo. Em cada
iteração, a aplicação das regras produz a configuração seguinte, que se
torna a base da próxima atualização.

\subsection{Parâmetros e restrições da entrada}\label{sec:entrada}

O arquivo informa $L$, $C$, $P$, a quantidade de threads $T$, a semente e
o limiar de ignição. Em seguida, especifica as componentes do vento e sua
intensidade, as quantidades de focos ($F$) e de zonas ($Z$), as coordenadas dos
focos e os retângulos de contenção com seus passos de ativação.

São exigidos $L,C,T,\mathrm{LIMIAR}>0$, $P,F,Z\geq0$,
componentes do vento em $\{-1,0,1\}$, direção diferente de $(0,0)$ e
intensidade entre 0 e 5. Os focos devem ser distintos, internos à grade e
posicionados sobre células combustíveis. As zonas devem ter limites
ordenados e internos à grade, com $0\leq p_{\mathrm{ativação}}<P$.
Consequentemente, não há passo de ativação válido quando $P=0$.

As duas aplicações recebem somente o caminho do arquivo de entrada:
\begin{lstlisting}
./fire_seq entrada.txt
./fire_omp entrada.txt
\end{lstlisting}
O valor de $T$ é obtido desse arquivo, inclusive nos experimentos que
variam a quantidade de threads.

\subsection{Cobertura, umidade e inicialização determinística}\label{sec:inicializacao}

A cobertura e a umidade são propriedades fixas de cada célula. A
inicialização percorre a matriz em ordem crescente de linhas e colunas,
com uma única thread. Para cada posição, primeiro obtém-se
\texttt{rand\_r(\&seed) \% 100} para determinar a cobertura e, imediatamente
depois, \texttt{rand\_r(\&seed) \% 101} para determinar a umidade.
Preservar essa ordem é necessário para que as versões recebam a mesma
configuração inicial no ambiente experimental.

\begin{table}[htbp]
\centering
\caption{Coberturas e parâmetros do modelo. As proporções são as da regra de geração, não contagens exatas de cada matriz.}
\label{tab:cobertura}
\begin{tabular}{clrrr}
\toprule
\textbf{Código} & \textbf{Cobertura} & \textbf{Sorteio} & \textbf{Fator} & \textbf{Queima} \\
\midrule
0 & Água & 0--9 & 0 & \pendente \\
1 & Solo exposto & 10--19 & 0 & \pendente \\
2 & Vegetação rasteira & 20--54 & 8 & 2 passos \\
3 & Floresta & 55--99 & 12 & 4 passos \\
\bottomrule
\end{tabular}
\end{table}

Água e solo exposto começam no estado não combustível. Vegetação rasteira
e floresta começam intactas, exceto nas posições dos focos iniciais, que
passam ao estado em chamas e recebem o tempo de queima de sua cobertura.
A quantidade de células inicialmente combustíveis é registrada antes da evolução,
incluindo os focos, para servir de denominador dos percentuais finais.

\subsection{Estados e estruturas de armazenamento}

O estado de uma célula pode ser não combustível (0), intacta (1), em
chamas (2), queimada (3) ou contenção (4). Cobertura e estado são conceitos
distintos: a cobertura identifica o material da célula; o estado indica
sua situação na evolução do incêndio.

\begin{table}[htbp]
\centering
\caption{Vetores utilizados para representar a floresta.}
\label{tab:vetores}
\begin{tabularx}{\linewidth}{lX}
\toprule
\textbf{Vetor} & \textbf{Função} \\
\midrule
\texttt{cobertura} & Tipo de cobertura de cada posição. \\
\texttt{umidade} & Umidade inteira entre 0 e 100. \\
\texttt{estado\_atual} & Configuração usada como base da atualização. \\
\texttt{proximo\_estado} & Configuração produzida pelo passo. \\
\texttt{tempo\_atual} & Tempo restante de queima na configuração atual. \\
\texttt{proximo\_tempo} & Tempo restante calculado para a configuração seguinte. \\
\texttt{ativacao} & Primeiro passo de contenção da posição, ou $-1$. \\
\bottomrule
\end{tabularx}
\end{table}

Os dois pares de vetores de estado e tempo evitam que uma célula utilize
uma atualização parcial de sua vizinha. Ao fim de cada passo, os
ponteiros de cada par são trocados. O armazenamento da grade é $O(LC)$; somadas as
descrições dos focos e zonas, o espaço é $O(LC+F+Z)$.

\subsection{Zonas de contenção}

Cada zona corresponde a um retângulo de limites inclusivos e a um passo
de ativação. Antes da simulação, constrói-se um mapa com o menor passo
associado a cada célula nas zonas sobrepostas; posições sem zona recebem
$-1$. A construção desse mapa não integra o trecho cronometrado.

No início do passo programado, apenas células intactas se tornam contenção.
Os demais estados permanecem inalterados: uma zona não apaga fogo existente
nem recupera células queimadas. A ativação ocorre somente no passo indicado,
e as células que se tornam contenção permanecem nesse estado.

\subsection{Vizinhança e potencial de ignição}

A propagação utiliza os oito vizinhos de Moore, ignorando posições fora
da grade. Apenas vizinhos em chamas contribuem. Para uma célula-alvo $(l,c)$
e um vizinho $(l_v,c_v)$, define-se a direção do vizinho para o alvo:
\begin{equation}
\Delta l=l-l_v,\qquad \Delta c=c-c_v.
\end{equation}
O peso básico $b_v$ vale 10 se $|\Delta l|+|\Delta c|=1$ e 7 para
vizinhos diagonais. Dadas a direção fixa $(w_l,w_c)$ e a intensidade $V$
do vento, o alinhamento e o peso ajustado são
\begin{align}
A_v &= \Delta l\,w_l+\Delta c\,w_c, \\
P_v &= \max(1,b_v+V A_v).
\end{align}
O produto escalar pode assumir valores entre $-2$ e $2$. Valores positivos
favorecem a propagação; negativos a desfavorecem. Quando $V=0$, o peso
ajustado coincide com o básico.

Se $f_i$ é o fator combustível e $u_i$ a umidade da célula-alvo, seu
potencial é
\begin{equation}
S_i=\sum_{v\text{ em chamas}}P_v,
\qquad I_i=\left\lfloor\frac{S_i f_i(100-u_i)}{100}\right\rfloor.
\end{equation}
Uma célula intacta se incendeia quando $I_i\geq\mathrm{LIMIAR}$.
Todos esses cálculos usam aritmética inteira, evitando diferenças de
arredondamento na propagação entre as versões.

\subsection{Transições, atualização síncrona e parada}\label{sec:transicoes}

Após a ativação das zonas, células não combustíveis, queimadas e de
contenção mantêm seus estados. Células intactas permanecem intactas ou
se incendeiam conforme o potencial. Células já em chamas têm seu tempo
reduzido em uma unidade e tornam-se queimadas ao atingir zero.

Uma célula recém-incendiada recebe o tempo completo de sua cobertura.
Ela não propaga fogo imediatamente dentro do mesmo passo: sua contribuição
passa a existir quando a configuração produzida se torna a configuração
atual. Essa separação preserva a atualização síncrona e permite que as
células sejam processadas em ordens diferentes sem alterar o resultado.

A sequência definida para cada passo é: ativar as zonas, produzir os
próximos estados e tempos, calcular as estatísticas do próximo estado,
trocar os vetores e verificar a parada. A execução termina ao completar
$P$ passos ou quando não há células em chamas. Se não houver fogo após a
inicialização, nenhum passo é executado, mesmo com contenções futuras.

\subsection{Estatísticas e saída}

As novas ignições de um passo são as transições de intacta para em chamas;
os focos iniciais não integram esse total. O pico registra a maior
quantidade de novas ignições e, em empate, o primeiro passo. Na ausência
de novas ignições, o resultado do pico é $(-1,0)$.

Se $B_0$ é a quantidade de células inicialmente combustíveis, $Q$ a quantidade final
queimada, $H$ a quantidade final em chamas e $K$ a quantidade final em
contenção, os percentuais são
\begin{equation}
\text{Queimado}=100\frac{Q+H}{B_0},\qquad
\text{Protegido}=100\frac{K}{B_0}.
\end{equation}
Quando $B_0=0$, ambos são zero. Células ainda em chamas entram no percentual
queimado porque já foram atingidas pelo incêndio.

O checksum é calculado sequencialmente após a simulação, na ordem linear
da matriz. Partindo de zero, para cada célula incorpora-se primeiro o estado
e depois o tempo, pela operação $h\leftarrow31h+x$, utilizando
\texttt{unsigned long long} conforme o enunciado. Ele complementa a
comparação dos resultados, sem substituir a verificação das demais saídas.

A saída contém o número de passos executados, as cinco contagens finais,
o total e o pico de ignições, os dois percentuais, o checksum e o tempo. Os percentuais têm duas
casas decimais e o tempo tem seis; todos os campos, exceto o tempo, devem
coincidir entre as versões.

\section{Implementação sequencial}

A versão sequencial, em \texttt{fire\_seq.c}, é a referência de correção e
de desempenho do trabalho. Ela segue diretamente a modelagem da
Seção~\ref{sec:modelagem}, sem otimizações que alterem a ordem das operações, de modo que
cada etapa do passo corresponde a uma função com responsabilidade única.

\subsection{Organização do código}

Os parâmetros de entrada, os focos, as zonas e os vetores da
Tabela~\ref{tab:vetores} são agrupados na estrutura \texttt{Simulacao},
passada por ponteiro a todas as funções. Focos e zonas são representados
pelas estruturas \texttt{Foco} e \texttt{ZonaContencao}. A
Tabela~\ref{tab:funcoes-seq} resume as funções e indica quais pertencem
ao trecho cronometrado.

\begin{table}[htbp!]
\centering
\caption{Funções da versão sequencial.}
\label{tab:funcoes-seq}
\begin{tabularx}{\linewidth}{lXc}
\toprule
\textbf{Função} & \textbf{Responsabilidade} & \textbf{Cronometrada} \\
\midrule
\texttt{ler\_entrada} & Leitura e validação do arquivo e dos limites de tamanho. & Não \\
\texttt{alocar\_estruturas} & Alocação dos sete vetores de tamanho $LC$. & Não \\
\texttt{gerar\_floresta} & Sorteio da cobertura e da umidade; estado inicial. & Não \\
\texttt{aplicar\_focos} & Marcação dos focos e validação da cobertura. & Não \\
\texttt{construir\_mapa\_ativacao} & Menor passo de ativação por célula. & Não \\
\texttt{ativar\_zonas} & Conversão de intactas em contenção no passo. & Sim \\
\texttt{potencial\_ignicao} & Cálculo de $I_i$ para uma célula intacta. & Sim \\
\texttt{calcular\_proximo\_estado} & Transições e contagem de ignições. & Sim \\
\texttt{trocar\_matrizes} & Troca dos ponteiros atual/próximo. & Sim \\
\texttt{existe\_em\_chamas} & Verificação da condição de parada. & Sim \\
\texttt{calcular\_checksum} & Checksum na ordem linear. & Não \\
\bottomrule
\end{tabularx}
\end{table}

\subsection{Preparação da simulação}

A leitura verifica o retorno de cada \texttt{fscanf} e aplica as
restrições da Seção~\ref{sec:entrada}, encerrando o programa com \texttt{EXIT\_FAILURE}
diante de qualquer inconsistência. A verificação de que os focos estão
sobre células combustíveis só é possível após a geração da floresta e,
por isso, é feita em \texttt{aplicar\_focos}. Como todos os índices da
simulação são do tipo \texttt{int}, a leitura também rejeita grades em
que $LC$ excede \texttt{INT\_MAX} ou em que o tamanho em bytes dos
vetores não cabe em \texttt{size\_t}; verificação análoga é aplicada às
quantidades de focos e zonas. Todas as alocações têm o retorno conferido.

Na geração da floresta, a semente lida é copiada para uma variável
local do tipo \texttt{unsigned int}, tipo exigido por \texttt{rand\_r},
evitando a conversão de ponteiros entre \texttt{int} e
\texttt{unsigned int}. As chamadas passam pela função
\texttt{fire\_rand\_r}, que utiliza o \texttt{rand\_r} nativo no
Linux e, no Windows e no macOS, uma reimplementação do algoritmo da
glibc. No Windows, \texttt{rand\_r} não existe; no macOS, existe, mas
emprega outro algoritmo e produziria uma floresta diferente para a mesma
semente. Assim, a configuração inicial é a mesma nas três plataformas. A ordem das chamadas é a descrita
na Seção~\ref{sec:inicializacao}: primeiro a cobertura e, em seguida, a
umidade da mesma célula.

O mapa de ativação é inicializado com $-1$ e preenchido percorrendo cada
retângulo; em zonas sobrepostas, prevalece o menor passo. Por fim, a
quantidade de células inicialmente combustíveis $B_0$ é contada a partir
da cobertura, antes do início da cronometragem.

\subsection{Laço principal}

O trecho cronometrado é delimitado por duas chamadas a
\texttt{omp\_get\_wtime()}; o cabeçalho \texttt{omp.h} é incluído
apenas para essa medição, e o programa não contém diretivas OpenMP.
O laço segue a ordem definida na Seção~\ref{sec:transicoes}:

\begin{lstlisting}
se existe_em_chamas():
    para passo = 0 ate P-1:
        ativar_zonas(passo)
        ignicoes = calcular_proximo_estado()
        total += ignicoes
        se ignicoes > pico_qtd: pico = (passo, ignicoes)
        trocar_matrizes()
        se nao existe_em_chamas(): interromper
\end{lstlisting}

A função \texttt{ativar\_zonas} percorre toda a grade e altera
\texttt{estado\_atual} diretamente, o que é seguro porque ocorre antes
da propagação do mesmo passo. Em \texttt{calcular\_proximo\_estado},
cada célula é tratada conforme seu estado: não combustíveis, queimadas e
contenções são copiadas com tempo zero; células em chamas têm o tempo
decrementado e passam a queimadas ao atingir zero; células intactas
chamam \texttt{potencial\_ignicao}. Todas as leituras são feitas em
\texttt{estado\_atual} e \texttt{tempo\_atual}, e todas as escritas em
\texttt{proximo\_estado} e \texttt{proximo\_tempo}, o que garante a
atualização síncrona.

O potencial percorre os deslocamentos $(d_l,d_c)\in\{-1,0,1\}^2$,
exceto $(0,0)$, e descarta vizinhos fora da grade. Para cada vizinho em
chamas, a direção até o alvo é $(-d_l,-d_c)$, a partir da qual se obtêm o
peso básico, o alinhamento $A_v$ e o peso $P_v$ com limite inferior 1.
O resultado final usa divisão inteira, equivalente ao piso para valores
não negativos.

O pico de ignições é atualizado apenas quando a contagem do passo é
estritamente maior que o máximo anterior, o que mantém o primeiro passo
em caso de empate. Como \texttt{pico\_passo} começa em $-1$ e
\texttt{pico\_qtd} em $0$, o resultado $(-1,0)$ é preservado quando não
há novas ignições. A troca das configurações é feita apenas pela troca
dos ponteiros, sem cópia de dados.

\subsection{Finalização e custo}

Após a cronometragem, são contados os estados finais, calculados os
percentuais com $B_0$ como denominador e obtido o checksum, acumulado em
\texttt{unsigned long long}, cujo estouro é aritmética modular bem
definida em C. A memória é liberada ao final.

Cada passo realiza três varreduras da grade: ativação das zonas, cálculo
do próximo estado e verificação de parada, esta última interrompida no
primeiro foco encontrado. Como cada célula intacta examina no máximo
oito vizinhos, o custo por passo é $O(LC)$ e o custo total do trecho
cronometrado é $O(P\,LC)$. O trabalho por célula, porém, não é uniforme:
células intactas são mais caras que as demais, o que é relevante para a
escolha do escalonamento na versão paralela.

\section{Implementação paralela}

A versão paralela, em \texttt{fire\_omp.c}, parte do código sequencial e
mantém idênticas a leitura, a validação, a alocação, a geração da
floresta, a aplicação dos focos, a construção do mapa de ativação e o
cálculo dos resultados finais. Todas essas etapas estão fora do trecho
cronometrado, e a geração da floresta deve, por exigência do enunciado,
ser feita por uma única thread. A paralelização concentra-se, portanto,
no laço de passos, encapsulado na função \texttt{simular}.

\subsection{Estratégia de paralelização}

Como a atualização é síncrona, o próximo estado de cada célula depende
apenas de \texttt{estado\_atual}, \texttt{tempo\_atual} e dos dados fixos,
e cada iteração escreve somente na própria posição de
\texttt{proximo\_estado} e \texttt{proximo\_tempo}. Não há, assim,
dependência entre iterações de um mesmo passo, e o laço sobre as células
pode ser dividido entre threads por decomposição de domínio sobre o
índice linear. Entre passos, por outro lado, a dependência é total: o
passo $p+1$ só pode começar depois que toda a configuração do passo $p$
estiver pronta.

Em vez de abrir uma região paralela a cada passo, o programa utiliza uma
única região \texttt{parallel} persistente, que envolve todo o laço
\texttt{while}. As threads são criadas uma vez e percorrem juntas os
passos, o que evita o custo repetido de criação e sincronização da
equipe ao longo de até $P$ passos. Dentro dela, cada passo é organizado
da seguinte forma:

\begin{lstlisting}
#pragma omp parallel num_threads(T)
while (continua):
    #pragma omp for schedule(static)        // ativar_zonas
    #pragma omp for schedule(guided, 1024) \
            reduction(+:ignicoes) reduction(||:tem_chamas)
    #pragma omp single                      // estatisticas, troca, parada
\end{lstlisting}

A ativação das zonas é feita em \texttt{ativar\_zonas} por um
\texttt{omp for} órfão, isto é, uma diretiva de compartilhamento de
trabalho fora do escopo léxico da região paralela, que se associa à
equipe em execução. Cada iteração lê \texttt{ativacao[i]}, que é fixo, e
escreve apenas em \texttt{estado\_atual[i]}, sem conflito entre
threads.

A atualização das células foi reorganizada na função
\texttt{atualizar\_celula}, que processa uma única posição e não contém
diretivas OpenMP; o laço que a chama é o que recebe o \texttt{omp for}.
Além do indicador de nova ignição, a função informa se a célula estará
em chamas no próximo estado. Esses valores são combinados por duas
reduções: soma para as novas ignições e disjunção lógica para a
existência de fogo. Com isso, a condição de parada passa a ser obtida
no próprio laço de atualização, e a varredura adicional de
\texttt{existe\_em\_chamas} a cada passo, presente na versão sequencial,
deixa de ser necessária. A função continua sendo usada apenas na
verificação inicial, antes da região paralela; se não houver fogo ou se
$P=0$, a região nem chega a ser aberta.

Em \texttt{potencial\_ignicao}, os oito deslocamentos de Moore são
armazenados em vetores constantes, e o laço sobre eles foi reescrito sem
desvios (\texttt{continue}), para permitir a vetorização com
\texttt{omp simd reduction(+:S)}. Um vizinho fora da grade tem o índice
substituído pelo da própria célula, sempre válido, e sua contribuição é
anulada pelo indicador \texttt{dentro}; o mesmo ocorre com vizinhos que
não estão em chamas. A soma $S_i$ resultante é a mesma da versão
sequencial.

\subsection{Compartilhamento e sincronização}

A região paralela usa \texttt{default(none)}, o que obriga a declarar
explicitamente o escopo de cada variável. São compartilhados o ponteiro
para a estrutura \texttt{Simulacao}, o tamanho $n$ e as variáveis de
controle do laço: \texttt{passo}, \texttt{continua}, o total e o pico
de ignições e os dois acumuladores das reduções. O índice dos laços
\texttt{omp for} e as variáveis locais de \texttt{atualizar\_celula} e
\texttt{potencial\_ignicao} são privados de cada thread. Os vetores
da simulação são compartilhados, mas nenhuma posição é escrita por mais
de uma thread no mesmo laço.

A corretude depende das barreiras implícitas ao final de cada construção
de compartilhamento de trabalho, que reproduzem a ordem da
Seção~\ref{sec:transicoes}:

\begin{itemize}
\item a barreira ao final da ativação garante que todas as contenções do
passo estejam aplicadas antes que qualquer thread avalie um vizinho;
\item a barreira ao final da atualização garante que todos os próximos
estados estejam escritos e que as reduções estejam combinadas antes do
bloco \texttt{single};
\item a barreira ao final do \texttt{single} garante que a troca de
ponteiros e os novos valores de \texttt{passo} e \texttt{continua} sejam
vistos por todas as threads antes da próxima iteração do
\texttt{while}.
\end{itemize}

As operações de custo constante do passo, como a atualização do total e
do pico, a troca dos vetores, o incremento de \texttt{passo} e o cálculo
de \texttt{continua}, são feitas por uma única thread no bloco
\texttt{single}. O mesmo bloco zera os acumuladores das reduções para o
passo seguinte, o que dispensa um segundo \texttt{single} e uma barreira
adicional por passo. Não há uso de \texttt{critical} nem de
\texttt{atomic}: toda a combinação de resultados parciais é feita pelas
reduções.

O número de threads é definido por \texttt{num\_threads(T)}, com o valor
lido do arquivo de entrada, e \texttt{omp\_set\_dynamic(0)} é chamado
antes da região para impedir que o ambiente reduza o tamanho da equipe.
Como as reduções operam sobre inteiros, cuja soma é associativa, os
resultados não dependem da ordem em que as threads combinam seus valores
parciais. As saídas são, portanto, as mesmas para qualquer quantidade de
threads e qualquer política de escalonamento.

\subsection{Escalonamento}

Os dois laços paralelos têm perfis de custo distintos e, por isso,
utilizam políticas diferentes. A ativação das zonas executa trabalho
praticamente uniforme por iteração: uma comparação e, eventualmente,
uma escrita. Para ela, usa-se \texttt{schedule(static)}, que atribui a
cada thread um bloco contíguo e fixo de índices, com sobrecarga mínima
e boa localidade de memória. A política é declarada explicitamente
porque, sem a cláusula, o escalonamento adotado depende da
implementação.

Na atualização das células, o custo por iteração é irregular. Células
intactas calculam o potencial sobre oito vizinhos, enquanto as demais
apenas copiam ou decrementam valores. A distribuição desse custo muda ao
longo da simulação, à medida que a frente de fogo avança e as regiões
queimadas e contidas crescem. Uma divisão estática em blocos iguais pode
atribuir a uma thread uma faixa com muito mais células intactas que a
de outra. Por isso, o laço usa \texttt{schedule(guided, 1024)}: os
blocos começam grandes e diminuem conforme o trabalho restante,
equilibrando a carga no final do laço. O tamanho mínimo de 1024
iterações limita a quantidade de requisições ao escalonador e preserva
a localidade, evitando a sobrecarga de políticas como
\texttt{dynamic} com blocos unitários.

A comparação entre esta política e outras alternativas, exigida pelo
enunciado, é apresentada na Seção~\ref{sec:resultados}.

\section{Validação da correção}


\begin{table}[htbp]
\centering
\caption{Validação da versão final (macOS).}
\begin{tabularx}{\linewidth}{@{}p{2.6cm}p{1.4cm}p{2.2cm}X@{}}
\toprule
\textbf{Caso} & \textbf{Threads} & \textbf{Schedule} & \textbf{Resultado} \\
\midrule
Carga pequena & 1--8 & guided(1024) & Checksum e passos idênticos entre \texttt{fire\_seq} e \texttt{fire\_omp} nas 90 execuções (checksum $17073017104646724864$; 80 passos) \\
Carga média & 1--8 & guided(1024) & Checksum e passos idênticos nas 90 execuções (checksum $17187382406529706254$; 100 passos) \\
Carga grande & 1--8 & guided(1024) & Checksum e passos idênticos nas 90 execuções (checksum $14825089465781038537$; 100 passos) \\
\bottomrule
\end{tabularx}
\end{table}

Para cada carga, o checksum e a quantidade de passos executados
coincidiram em todas as 90 execuções (versão sequencial e versão
OpenMP com $T$ de 2 a 8, dez repetições cada), confirmando que o
resultado independe da quantidade de threads, conforme exigido pelo
enunciado. Nas três cargas a simulação consumiu todos os $P$ passos
configurados, isto é, o incêndio não se extinguiu antes do limite
definido na Tabela de cargas. As entradas com zero passos, sem focos e
com zonas de contenção sobrepostas foram contempladas na validação da
entrada (Seção~\ref{sec:entrada}) e na lógica de \texttt{construir\_mapa\_ativacao},
mas ainda não possuem uma execução automatizada dedicada nos scripts de
teste; essa cobertura será complementada nos experimentos do cluster
Linux.


\section{Ambiente e metodologia experimental}\label{sec:ambiente}

\subsection{Ambiente das duas versões}
\begin{table}[htbp]
\centering
\caption{Ambiente experimental --- macOS.}
\begin{tabularx}{\linewidth}{lX}
\toprule
\textbf{Item} & \textbf{Configuração} \\
\midrule
CPU e núcleos & Apple M4 (4 núcleos de desempenho + 6 de eficiência; 10 núcleos no total, sem SMT/hyperthreading) \\
Memória e sistema & 16\,GiB de RAM; macOS 26.5.2 (build 25F84); arquitetura arm64 (Apple Silicon) \\
Compilador & Apple clang 17.0.0 (clang-1700.6.4.2), \texttt{target=arm64-apple-darwin25.5.0}, com \texttt{libomp} do Homebrew (\texttt{/opt/homebrew/opt/libomp}) \\
Compilação sequencial & \texttt{Makefile.macos}, alvo \texttt{fire\_seq}: \texttt{clang -O2 -Wall -Wextra -std=c11 -Xpreprocessor -fopenmp -I<libomp>/include -L<libomp>/lib -lomp fire\_seq.c -o fire\_seq} \\
Compilação paralela & mesmos parâmetros de compilação, alvo \texttt{fire\_omp} (\texttt{fire\_omp.c}) \\
Execução OpenMP & \texttt{T} lido do arquivo de entrada e fixado com \texttt{num\_threads(T)} e \texttt{omp\_set\_dynamic(0)}; a política de escalonamento do laço de atualização (\texttt{guided(1024)} nesta rodada) é fixada em tempo de compilação na diretiva \texttt{omp for} \\
Versão dos fontes & commit \texttt{fe355af} (branch \texttt{main}) \\
\bottomrule
\end{tabularx}
\end{table}

\noindent\textbf{Observação sobre o ambiente.} Todos os resultados
apresentados na Seção~\ref{sec:resultados} e nas seções seguintes desta
versão do relatório foram obtidos exclusivamente no ambiente macOS
descrito acima, com o script \texttt{scripts/run\_macos.sh}. Os
experimentos equivalentes no cluster Linux da disciplina (usando
\texttt{scripts/run\_linux.sh} e \texttt{Makefile}) ainda serão
executados; quando disponíveis, serão incorporados a uma versão
revisada deste relatório, junto com a comparação de uma segunda
política de escalonamento exigida pelo enunciado.

\subsection{Cargas e protocolo de execução}
\begin{table}[htbp]
\centering
\caption{Cargas fornecidas. $T$ indica o valor original, que pode variar nos experimentos.}
\begin{tabular}{lrrrrrr}
\toprule
\textbf{Carga} & $L$ & $C$ & $P$ & $T$ & $F$ & $Z$ \\
\midrule
Pequena & 400 & 500 & 80 & 4 & 3 & 2 \\
Média & 1200 & 1500 & 100 & 8 & 6 & 4 \\
Grande & 2500 & 2500 & 100 & 8 & 10 & 6 \\
\bottomrule
\end{tabular}
\end{table}

Cada configuração foi executada pelo script \texttt{scripts/run\_macos.sh}:
a versão sequencial é executada uma vez por repetição (sem variar
threads) e a versão paralela é executada para $T\in\{1,\dots,8\}$, com o
$4^{\text{o}}$ campo da primeira linha da entrada sobrescrito
temporariamente a cada valor de $T$ testado; os demais parâmetros da
Tabela acima são mantidos. Foram realizadas 10 repetições por combinação
de carga e $T$, totalizando $3\times(10+8\times10)=270$ execuções nesta
rodada. Nesta versão do relatório, apenas a política
\texttt{guided(1024)} foi medida no macOS; a comparação com uma segunda
política de escalonamento, exigida pelo enunciado, será obtida junto com
os experimentos no cluster Linux.

\subsection{Delimitação do trecho cronometrado}

O recorte especificado inclui a ativação das zonas, a propagação, a
atualização dos estados e tempos, as estatísticas, a troca das matrizes e a
verificação da parada. Exclui a leitura, a validação, a alocação, a geração
da floresta, a definição dos focos, a construção do mapa, o checksum, os
percentuais e a impressão. A medição
deve utilizar \texttt{omp\_get\_wtime()}.


\subsection{Métricas}
Para a mesma carga, o speedup e a eficiência percentual são definidos por
\begin{equation}
S(T)=\frac{t_{\mathrm{seq}}}{t_{\mathrm{omp}}(T)},\qquad
E(T)=\frac{S(T)}{T}\times100\%.
\end{equation}
A referência é o programa sequencial. A versão OpenMP com uma thread é
uma configuração adicional e não substitui essa referência.

\section{Resultados}\label{sec:resultados}

\subsection{MacOs}
Os resultados desta seção foram obtidos no ambiente macOS descrito na
Seção~\ref{sec:ambiente}, com a política \texttt{guided(1024)}. Cada
valor de $t_{\mathrm{seq}}$ e $t_{\mathrm{omp}}$ é a média aritmética de
10 repetições; o script de execução também registrou o desvio entre
repetições (não incluído na tabela), pequeno para as cargas média e
grande e um pouco maior, em termos relativos, para a carga pequena, cujo
tempo de execução é da ordem de dezenas de milissegundos. Os
experimentos no cluster Linux da disciplina, incluindo a comparação com
uma segunda política de escalonamento, serão incorporados a uma versão
revisada desta tabela e dos gráficos.

\begin{table}[htbp]
\centering\small
\caption{Tempos e métricas de desempenho (macOS, Apple M4, política \texttt{guided(1024)}).}
\begin{tabular}{llrrrrr}
\toprule
\textbf{Carga} & \textbf{Schedule} & $T$ & $t_{\mathrm{seq}}$ (s) & $t_{\mathrm{omp}}$ (s) & $S(T)$ & $E(T)$ (\%) \\
\midrule
\textbf{Pequena}  & guided(1024) & 2 & 0.083092 & 0.052361 & 1.59 & 79.35 \\
                  & guided(1024) & 3 & 0.083092 & 0.035459 & 2.34 & 78.11 \\
                  & guided(1024) & 4 & 0.083092 & 0.032095 & 2.59 & 64.72 \\
                  & guided(1024) & 5 & 0.083092 & 0.029900 & 2.78 & 55.58 \\
                  & guided(1024) & 6 & 0.083092 & 0.028959 & 2.87 & 47.82 \\
                  & guided(1024) & 7 & 0.083092 & 0.027257 & 3.05 & 43.55 \\
                  & guided(1024) & 8 & 0.083092 & 0.033282 & 2.50 & 31.21 \\
\midrule
\textbf{Média}    & guided(1024) & 2 & 0.942207 & 0.588427 & 1.60 & 80.06 \\
                  & guided(1024) & 3 & 0.942207 & 0.456809 & 2.06 & 68.75 \\
                  & guided(1024) & 4 & 0.942207 & 0.376859 & 2.50 & 62.50 \\
                  & guided(1024) & 5 & 0.942207 & 0.346127 & 2.72 & 54.44 \\
                  & guided(1024) & 6 & 0.942207 & 0.328233 & 2.87 & 47.84 \\
                  & guided(1024) & 7 & 0.942207 & 0.300556 & 3.14 & 44.78 \\
                  & guided(1024) & 8 & 0.942207 & 0.278883 & 3.38 & 42.23 \\
\midrule
\textbf{Grande}   & guided(1024) & 2 & 3.302497 & 2.267915 & 1.46 & 72.81 \\
                  & guided(1024) & 3 & 3.302497 & 1.727732 & 1.91 & 63.72 \\
                  & guided(1024) & 4 & 3.302497 & 1.422973 & 2.32 & 58.02 \\
                  & guided(1024) & 5 & 3.302497 & 1.341581 & 2.46 & 49.23 \\
                  & guided(1024) & 6 & 3.302497 & 1.258350 & 2.62 & 43.74 \\
                  & guided(1024) & 7 & 3.302497 & 1.252327 & 2.64 & 37.67 \\
                  & guided(1024) & 8 & 3.302497 & 1.181848 & 2.79 & 34.93 \\
\bottomrule
\end{tabular}
\end{table}

\grafico{figuras/tempos.pdf}{Tempo em função do número de threads. Identificar cargas e schedules.}{fig:tempos}
\grafico{figuras/speedup.pdf}{Speedup em função do número de threads, com referência ideal $S(T)=T$.}{fig:speedup}
\grafico{figuras/eficiencia.pdf}{Eficiência percentual em função do número de threads.}{fig:eficiencia}

\section{Análise e discussão}

\textbf{Todas as observações abaixo se referem exclusivamente ao
ambiente macOS} descrito na Seção~\ref{sec:ambiente} e à política
\texttt{guided(1024)}; conclusões sobre outros escalonamentos e sobre o
comportamento no cluster Linux ficam para a versão revisada deste
relatório.

\textbf{Ganho de threads limitado pelo número de núcleos de desempenho.}
O Apple M4 utilizado possui apenas 4 núcleos de desempenho e 6 de
eficiência, sem SMT. Para as três cargas, o speedup cresce de forma
quase linear até $T=3$--$4$ (por exemplo, $S(4)=2.50$ na carga média) e
passa a crescer bem mais devagar a partir daí: entre $T=4$ e $T=8$ a
carga média ganha apenas $0.88$ de speedup adicional, contra $1.66$
entre $T=1$ e $T=4$. Esse comportamento é compatível com threads
adicionais sendo atendidas por núcleos de eficiência, mais lentos que os
de desempenho, e explica por que $E(T)$ cai de forma monótona em todas
as cargas: de $83$--$84\%$ em $T=1$ para $31$--$43\%$ em $T=8$.

\textbf{A carga grande escala pior do que a carga média.} Em $T=8$, a
carga média atinge $S(8)=3.38$ e $E(8)=42.23\%$, enquanto a carga grande
atinge apenas $S(8)=2.79$ e $E(8)=34.93\%$, apesar de ter mais trabalho
por passo ($2500\times2500$ contra $1200\times1500$ células). Esse
resultado é o oposto do que a lei de Amdahl prevista isoladamente
sugeriria (cargas maiores tendem a diluir melhor a fração serial fixa
por passo). Uma explicação provável é a pressão sobre a hierarquia de
memória: a grade da carga grande ($\approx 6{,}25$ milhões de células,
com sete vetores de tamanho $LC$) não cabe nos caches compartilhados
pelos núcleos de desempenho, de modo que, ao aumentar o número de
threads ativas, a carga passa a ser limitada pela largura de banda de
memória, e não apenas pelo paralelismo de cálculo. Essa hipótese deve
ser reexaminada no cluster Linux, cuja hierarquia de cache e memória é
distinta.

\textbf{Queda de speedup na carga pequena em $T=8$.} A carga pequena tem
o pior comportamento relativo em threads altas: o speedup sobe até
$S(7)=3.05$ e cai para $S(8)=2.50$, o único caso, entre as três cargas,
em que aumentar $T$ piora o tempo de execução. Como o trecho
cronometrado da carga pequena dura poucas dezenas de milissegundos e
tem apenas $200\,000$ células distribuídas por passo, o custo fixo de
cada barreira implícita e de cada região de trabalho compartilhado
(Seção~\ref{sec:transicoes}) passa a ser proporcionalmente maior à
medida que $T$ cresce e que a fatia de trabalho por thread diminui;
com 8 threads, cada uma recebe em média $25\,000$ células por passo,
valor próximo do tamanho mínimo de bloco do \texttt{guided(1024)}, o
que reduz a eficácia do balanceamento de carga e evidencia o custo de
sincronização por passo. Esse efeito é uma evidência de que, para
cargas pequenas, threads demais podem ser contraproducentes com essa
política de escalonamento --- um ponto a ser confrontado com a segunda
política de escalonamento ainda pendente.

\textbf{Consistência dos resultados.} Como reportado na
Seção~\ref{sec:resultados} anterior (validação), o checksum e a
quantidade de passos executados foram idênticos entre \texttt{fire\_seq}
e \texttt{fire\_omp} para todo $T$ de 1 a 8, nas três cargas, confirmando
que a paralelização não altera o resultado da simulação --- apenas o seu
tempo de execução, exatamente como esperado da modelagem com atualização
síncrona e reduções (Seção~\ref{sec:transicoes}).

\section{Conclusão}

As implementações sequencial e paralela atendem aos requisitos do
enunciado: mesma saída determinística para uma dada entrada,
independentemente da quantidade de threads, com a versão paralela
construída sobre uma região \texttt{parallel} persistente, \texttt{omp
for} com \texttt{default(none)}, reduções para os contadores e SIMD no
cálculo do potencial de ignição. No ambiente macOS avaliado (Apple M4,
10 núcleos, política \texttt{guided(1024)}), a versão OpenMP produziu
speedups entre $2.50\times$ e $3.38\times$ em relação à versão
sequencial, com eficiência decrescente à medida que a quantidade de
threads ultrapassa o número de núcleos de desempenho da máquina. A carga
grande escalou pior do que a média, indicando possível limitação de
banda de memória, e a carga pequena chegou a piorar ao usar 8 threads,
evidenciando que o custo de sincronização por passo pode superar o
ganho de paralelismo quando a carga de trabalho é pequena.

Esses resultados são específicos do ambiente macOS descrito na
Seção~\ref{sec:ambiente} e de uma única política de escalonamento. Os
próximos passos deste trabalho são: (i) executar o mesmo protocolo
experimental no cluster Linux da disciplina, com \texttt{scripts/run\_linux.sh}
e o \texttt{Makefile} padrão; (ii) medir pelo menos uma segunda política
de escalonamento (por exemplo, \texttt{static} ou \texttt{dynamic}) para
o laço de atualização, conforme exigido pelo enunciado; e (iii)
completar a validação de correção com os casos de zero passos/sem focos
e de zonas de contenção sobrepostas. Espera-se que o cluster Linux, com
mais núcleos e uma hierarquia de memória diferente da do M4, permita
avaliar se a queda de eficiência observada aqui é uma limitação da
arquitetura específica ou do algoritmo de paralelização em si.


\begin{thebibliography}{9}
\bibitem{enunciado}
Universidade de São Paulo, ICMC/SSC.
\textit{SSC0903 --- Primeiro Trabalho Prático: simulação paralela da
propagação direcional de incêndio com zonas de contenção}.
Arquivo \texttt{Trab01-Fire-OMP-v1.pdf}, setembro de 2026.
\end{thebibliography}

\end{document}
